\section{Proposed Method: AGS}
\label{sec:extension}
Agreement-Gated Search (AGS) allocates concept search according to the input. The component results show diminishing joint gains from search and visual prompting, alongside substantial variation across queries. The proposed hypothesis is that agreement between text-only and visual-only masks identifies inputs for which further search has limited value. AGS applies input-dependent computation allocation~\cite{branchynet,adaptiveinference} to the concept-search stage of CG-ICS.

\subsection{Decision Rule}
For the first reference-derived concept $c_0$, let $M_t(c_0)$ denote its text-only query mask and $M_v$ the visual-only mask. If either probe is empty, set $A=0$; otherwise define
\begin{equation}
    a=\operatorname{IoU}\!\left(M_t(c_0),M_v\right),\qquad
    A=\mathbf{1}\!\left[a\geq\tau\right].
\end{equation}
If $A=1$, the system retains $c_0$ and produces the joint text-plus-visual mask. Otherwise, it performs concept search before the final joint prediction:
\begin{equation}
    \widehat M_q=
    \begin{cases}
        M_{tv}(c_0), & A=1,\\
        M_{\mathrm{search}}, & A=0.
    \end{cases}
    \label{eq:gate}
\end{equation}
The primary threshold is fixed at $\tau=0.7$, with 0.5 and 0.9 used for sensitivity analysis. No query annotation enters this decision. Agreement measures consistency between the probe masks; its relationship to their target accuracy is formalized in Appendix~\ref{app:math}.

\subsection{Cost--Accuracy Analysis}
Let $A=1$ mean that the gate accepts the initial concept. After a shared
prefix of cost $C_p$, let $C_e$ and $C_s$ be the remaining costs of the
early and search branches, including their final segmentation calls.
The gated and staged always-search costs satisfy
\begin{equation}
  \begin{aligned}
  C_g&=C_p+A C_e+(1-A)C_s,\\
  C_{\mathrm{st}}&=C_p+C_s.
  \end{aligned}
  \label{eq:gate-cost}
\end{equation}
Writing $D=C_s-C_e$ and denoting the corresponding final-mask IoUs by
$J_e$ and $J_s$ gives
\begin{equation}
  \begin{aligned}
  \mathbb{E}[C_g-C_{\mathrm{st}}]&=-\mathbb{E}[A D],\\
  \mathbb{E}[J_g-J_{\mathrm{st}}]&=\mathbb{E}[A(J_e-J_s)].
  \end{aligned}
  \label{eq:gate-decomposition}
\end{equation}
The first identity measures computation avoided on accepted inputs;
the second measures the accuracy change on those same inputs. If
$H=C_{\mathrm{st}}-C_{\mathrm{full}}$ is the staging overhead relative to
original CG-ICS, the gate is faster in expectation precisely when
$\mathbb{E}[A D]>\mathbb{E}[H]$. Consequently, an acceptance rate alone
does not establish a speedup: the saved branch costs must exceed the
probe overhead. The staged always-search and random-routing controls
separate this overhead from the value of the agreement signal.



\subsection{Evaluation Design}
The evaluation design separates four questions. A fixed-candidate-pool ablation compares ranking by $s_rs_q$, $s_r$, and $s_q$ using the same candidate evidence and final visual prompts. A fresh-query comparison tests the original full system, the single-concept route, and the one-initial-candidate setting. Domain transfer uses SUIM underwater images~\cite{suim} and public ISIC~2018 Task~1 lesion images and masks~\cite{isic,ham10000}. Finally, the agreement gate is compared with original full search, staged always-search, an initial-score gate, and random routing.

The protocol specifies 400 fresh COCO and 1,000 fresh LVIS query images, each with up to three references, plus 1,000 queries each from SUIM and ISIC (Appendix~\ref{app:protocol}). The initial-score threshold is 0.8; random policies search with probabilities 0.25, 0.5, and 0.75. These controls test whether agreement identifies useful opportunities to stop, beyond the benefit of simply performing less search.

The evaluation criteria are paired accuracy, latency distributions, reference sensitivity, search and request counts, startup, and owned-process memory.
